\section{Conclusión}

La ontología del presente en sistemas complejos articulada en torno al marco del Reloj Extramental Discreto (RECD) caracteriza el presente como un fenómeno \textbf{estratificado, fragmentado y jerárquico por capas}, en el que la intensificación local de la persistencia (hiper-persistencia + trapping estructural) constituye una forma de mayor densidad y adhesión, pero cuya capacidad de generar reorganización estructural depende críticamente de la reducción de la fragmentación, regulada por la antisincronización.

La antisincronización funciona como regulador de un efecto estructural de filtrado probabilístico bidireccional que modula la probabilidad de que los presentes locales intensificados se integren en niveles superiores (los términos «resistencia ascendente» y «resistencia descendente» se usan solo ocasionalmente como etiquetas heurísticas secundarias). Cuando este efecto de filtrado probabilístico disminuye (es decir, cuando la antisincronización desciende), se crean las condiciones para que la intensificación local adquiera efectos globales. La «memoria activa», entendida como retención dinámica y no pasiva (hiper-persistencia + trapping estructural), proporciona el soporte necesario para que tales intensificaciones puedan tener peso ontológico.

El tiempo extramental no aparece como un flujo uniforme ni como una sucesión de instantes discretos carentes de contenido, sino como un tiempo articulado por transiciones entre diferentes niveles de densidad, adhesión y coherencia de la persistencia. Las capas del presente se modulan mutuamente mediante efectos de filtrado probabilístico bidireccional, y las transiciones entre ellas constituyen los eventos ontológicamente significativos.

El trabajo representa un estado intermedio de un programa de investigación en desarrollo. Las definiciones y relaciones establecidas están sujetas a refinamiento a medida que se acumule nueva evidencia empírica —especialmente en regímenes con acoplamiento asimétrico, en sistemas de mayor dimensionalidad y en contextos empíricos adicionales— y a medida que se profundice el diálogo con tradiciones filosóficas de la temporalidad.

Esta ontología del presente no pretende sustituir otras concepciones del tiempo, sino ofrecer un marco preciso y operacionalmente anclado en la dinámica de sistemas complejos que permita pensar la discreción temporal como una propiedad emergente de la organización jerárquica de la persistencia.